\documentclass{situationist-preprint}

\title{The Existence of Social Kinds}
\msid{MET-20-a}
\msfamily{MET-20}
\msdiscipline{Metaphysics}
\mstype{Principal manuscript}
\msstatus{Revised 7 October 2026}
\mskeywords{social ontology; social kinds; collective acceptance; institutions; persistence; mind-dependence}

\begin{document}
\maketitle

\begin{abstract}
Money, borders, universities and marriages are kinds whose instances exist only because of what people collectively do and think. Are these kinds real, or are they ways of talking about the behaviour of individuals? We argue that social kinds exist. They figure in explanations that cannot be reformulated without them, and they persist independently of the attitudes of any particular individual. We give an account of their persistence: a social kind, or an institution that instantiates one, persists as long as it is not collectively rejected by the community whose practices contain its rules. Active acceptance is required for an institution to come into existence, but not for it to continue. We argue that dependence on collective attitudes of this kind does not make social kinds less real than natural kinds.
\end{abstract}

\paragraph{Note on the revision.} The version of 6 October 2026 stated the persistence condition (Lemma~\ref{l:persist}) as continued collective acceptance. The argument of Section~3 requires, and the cases of Section~3.2 support, the weaker condition stated here: absence of collective rejection. The two conditions are not equivalent, since an institution that no one actively endorses but no one rejects persists under the second and not under the first. The main theorem is restated accordingly. Two manuscripts that relied on the earlier formulation, MET-20-b and MET-20-c, have been withdrawn, together with POL-07-b; see the history of the collection.

\section{The question}

Some kinds are social: whether something belongs to them depends on the attitudes and practices of people. A piece of paper is a banknote because it is issued by a central bank and accepted in exchange; a line on the ground is a border because states and their citizens treat it as one; a group of people is a committee because it has been constituted as one by a body with the authority to do so. The question of this manuscript is whether such kinds exist, in the sense in which natural kinds exist (MET-27), or whether talk of them is a convenient way of describing the behaviour of individuals.

Three positions dominate the literature. \emph{Eliminativists} hold that there are only individuals and their attitudes, and that social kinds are not part of the inventory of the world. \emph{Constructivists}, following \citet{Searle1995,Searle2010}, hold that social kinds exist and are created by collective intentionality, through declarations of the form \enquote{X counts as Y in context C}. \emph{Rule-based} accounts, following \citet{GualaHindriks2015} and \citet{Guala2016}, identify institutions with rules in equilibrium: solutions to coordination problems that people have reason to follow because others follow them. Other views locate the existence of social kinds in the grounding of social facts \citep{Epstein2015} or in the conferral of properties by those with standing to confer them \citep{Asta2018}.

We defend the claim, shared by constructivist and rule-based accounts, that social kinds exist, and we add an account of the condition under which they persist, which these accounts leave underdetermined.

\subsection{A short history}

The question has a long history in the social sciences before it became a question in metaphysics. \citet{Durkheim1895} held that social facts are \enquote{things}: ways of acting, thinking and feeling that are external to individuals and exercise constraint over them, and that must be explained by other social facts rather than by facts about individuals. Weber's methodological individualism took the opposite view, that social collectives are to be treated as the resultants of the actions of individual persons, and that sociological explanation must ultimately be couched in terms of individual action. The disagreement between the two traditions was methodological as much as ontological, and much of twentieth-century social theory can be read as an attempt to find an intermediate position.

Analytic philosophy entered the debate through the theory of convention. \citet{Lewis1969} showed how conventions could arise as solutions to coordination problems sustained by mutual expectations, without explicit agreement. \citet{Gilbert1989} argued that many social phenomena require a stronger notion, joint commitment, by which individuals constitute themselves as a plural subject. \citet{Searle1995} combined collective intentionality with the assignment of status functions, and \citet{Tuomela2013} developed a detailed account of group attitudes. More recently, \citet{GualaHindriks2015} have shown that Searle's status functions and Lewis's conventions can be unified: a status function is a rule in equilibrium, described in a vocabulary that makes its deontic consequences explicit. The present manuscript takes this unified account as its starting point and asks what it implies about persistence.

\section{Framework}

\begin{definition}\label{d:social}
A kind $K$ is \emph{social} if and only if, for every instance $x$ of $K$, $x$ is a $K$ at least partly in virtue of the attitudes or practices of a community.
\end{definition}

\begin{definition}\label{d:accept}
A community \emph{collectively accepts} a status function \enquote{X counts as Y in C} if its members, in the relevant contexts, treat things that satisfy X as having the status Y, and expect others to do so \citep{Searle1995}.
\end{definition}

\begin{definition}\label{d:reject}
A community \emph{collectively rejects} a status function if its members, in the relevant contexts, refuse to treat things that satisfy X as having the status Y, and expect others to refuse as well.
\end{definition}

\begin{lemma}\label{l:incompatible}
A community cannot both collectively accept and collectively reject the same status function at the same time, in the same contexts.
\end{lemma}

\begin{proof}
Collective acceptance requires that members treat things satisfying X as having status Y in the relevant contexts, and collective rejection that they refuse to do so in those contexts. A member cannot both treat and refuse to treat the same thing in the same context, and the expectations that each attitude involves are incompatible.
\end{proof}

It follows from Lemma~\ref{l:incompatible} that acceptance entails non-rejection, so that the persistence condition defended below is weaker than the acceptance condition: every institution that persists by the acceptance condition persists by ours. Collective acceptance and collective rejection are not, however, contradictories. A community may neither accept nor reject a status function: its members may have stopped treating things as having the status, without refusing to treat them so, as when an institution falls out of use without being abolished.

\begin{principle}[Explanatory criterion]\label{p:criterion}
A kind exists if it has instances and figures in true explanations that cannot be reformulated without loss of explanatory power in terms of other kinds.
\end{principle}

Principle~\ref{p:criterion} is a version of the criterion that underlies realism about natural kinds (MET-27) and about the posits of successful sciences (SCI-02). We use it as a sufficient condition only.

\section{The argument}

\subsection{Existence}

\begin{lemma}\label{l:explain}
Social kinds figure in explanations that cannot be reformulated without them.
\end{lemma}

\begin{proof}
Consider the explanation of inflation by an increase in the money supply. The explanation generalizes over money, whatever its physical realization: banknotes, coins, account balances. Reformulated in terms of the attitudes of individuals, it would have to quantify over the open-ended and heterogeneous set of attitudes that make different things money in different communities, and it would lose the generalization that the same increase in the money supply has the same effect regardless of how money is realized. The same holds for explanations of trade by borders, of career paths by universities, and of the distribution of property by marriage. These explanations are true, and the kinds they invoke cannot be eliminated without loss.
\end{proof}

\subsection{Persistence}

The second part of the argument concerns the conditions under which social kinds, and the institutions that instantiate them, continue to exist.

\begin{lemma}\label{l:persist}
An institution that has come into existence through collective acceptance persists as long as it is not collectively rejected by the community whose practices contain its rules.
\end{lemma}

\begin{proof}
Consider three kinds of case. First, laws that remain in force although no one applies them: a statute that has not been repealed remains law, and a court asked to apply it will do so, even if no one has thought of it for decades. Second, currencies that remain legal tender although they are no longer used in exchange: a denomination that has fallen out of use is still money in the jurisdiction that issued it, as is shown by the fact that a creditor cannot refuse it in payment of a debt. Third, offices that remain in existence although they are vacant and no one attends to them: an unfilled position in an organization persists until the organization abolishes it.

In each case, collective acceptance has lapsed: members of the community do not, in the relevant contexts, treat things as having the status, since the contexts no longer arise. Yet the institution persists, as is shown by what happens when the contexts arise again: the law is applied, the coin is accepted, the office is filled. What would end the institution is not the lapse of acceptance but its rejection: repeal of the law, demonetization of the currency, abolition of the office. The persistence condition is therefore the absence of collective rejection.

The qualification \enquote{by the community whose practices contain its rules} is needed for the case in which an entire community forgets an institution, not because it has rejected it but because its practices no longer contain any trace of it. An institution of a vanished society does not persist merely because no one has rejected it. Following \citet{Epstein2015}, we say that the rules of an institution must remain anchored in the practices of some community, for instance in its legal records or its customs, for the absence of rejection to sustain it.
\end{proof}

Two further classes of case confirm the lemma.

\paragraph{Desuetude.} Legal systems differ on whether a law can lapse through disuse. English law does not recognize a doctrine of desuetude: a statute remains in force until it is repealed, however long it has gone unenforced. Scots law recognizes such a doctrine for some older statutes, which may cease to have effect when they have long been disregarded and a contrary practice has grown up. The difference confirms Lemma~\ref{l:persist} rather than undermining it. Where desuetude is recognized, the legal system contains a rule according to which prolonged disuse together with contrary practice amounts to rejection; where it is not recognized, disuse alone does not. In both systems, what ends a law is the rejection of it, by repeal or by a practice the system treats as equivalent to repeal; in neither does the mere lapse of active acceptance end it.

\paragraph{Demonetization.} When a state withdraws a currency, it does not wait for people to stop using it. It declares that certain notes or coins will cease to be legal tender on a given date, and from that date they are no longer money, however many people continue to hold them and however willing some might be to accept them. In the Indian demonetization of 2016, for instance, high-denomination notes ceased to be legal tender by announcement, while they were still in wide use. Here an institution ends while acceptance continues, because it is rejected by those with standing to reject it. The case is the converse of the dormant law, and the two together show that rejection, not acceptance, is what the persistence of institutions tracks.

\begin{remark}
The version of Lemma~\ref{l:persist} published on 6 October 2026 required continued collective acceptance. The cases in the proof show that this condition is too strong: it would imply that dormant laws, unused currencies and vacant offices cease to exist, which is contrary to the practices that sustain them. The error affected the arguments of the withdrawn manuscripts MET-20-b, on borders, and MET-20-c, on the persistence of institutions, which inferred from the lapse of acceptance to the cessation of the institution, and POL-07-b, which used the same inference in an argument about the stabilization of states.
\end{remark}

\subsection{The main result}

\begin{theorem}\label{t:main}
Social kinds exist. Their instances depend for their existence on the absence of collective rejection, rather than on active acceptance, once they have come into existence.
\end{theorem}

\begin{proof}
By Lemma~\ref{l:explain} and Principle~\ref{p:criterion}, social kinds exist. By Definition~\ref{d:social}, their instances depend on collective attitudes; by Lemma~\ref{l:persist}, the attitude on which their continued existence depends is the absence of collective rejection.
\end{proof}

\begin{corollary}\label{c:real}
Dependence on collective attitudes does not make social kinds less real than natural kinds.
\end{corollary}

\begin{proof}
Principle~\ref{p:criterion} is the criterion by which natural kinds are counted as real, and social kinds satisfy it. That social kinds depend on attitudes is a fact about what grounds their instances, not about whether they exist. Response-dependent properties such as colours (MET-17) also depend on the responses of perceivers, and are no less real for that.
\end{proof}

\section{Three persistence conditions compared}

Lemma~\ref{l:persist} can be compared with the two alternatives found in the literature. Write $\mathrm{Acc}$ for continued collective acceptance, $\mathrm{Eq}$ for continued equilibrium of the institution's rule, in the sense of \citet{GualaHindriks2015}, and $\neg\mathrm{Rej}$ for the absence of collective rejection by an anchoring community.

\begin{center}\small
\begin{tabular}{@{}lccc@{}}
\toprule
Case & $\mathrm{Acc}$ & $\mathrm{Eq}$ & $\neg\mathrm{Rej}$ \\
\midrule
Dormant statute, still in force & no & no & yes \\
Unused denomination, still legal tender & no & no & yes \\
Vacant office, not abolished & no & no & yes \\
Demonetized note, still accepted & yes & yes & no \\
Repealed law, still observed by habit & yes & yes & no \\
Institution of a vanished society & no & no & no \\
\midrule
Verdict of practice & \multicolumn{3}{c}{persists, persists, persists, ends, ends, ends} \\
\bottomrule
\end{tabular}
\end{center}

\begin{proposition}\label{p:compare}
On the six cases, the absence of collective rejection by an anchoring community agrees with the verdict of practice in every case; continued acceptance and continued equilibrium each disagree with it in five.
\end{proposition}

\begin{proof}
By inspection of the table. Acceptance and equilibrium classify the first three cases as ending, against practice, and the fourth and fifth as persisting, against practice; they agree with practice only on the sixth.
\end{proof}

The comparison should not be overstated. The cases were chosen to separate the conditions, and a defender of the acceptance condition may reply that the dormant cases involve acceptance by officials, who would apply the law or accept the coin if asked. This reply, however, concedes the point: what the officials have is not active acceptance but a disposition not to reject, which is the condition Lemma~\ref{l:persist} states. The equilibrium account faces a related difficulty. A rule that is not currently followed, because the situations it governs do not arise, is not in equilibrium in any behavioural sense, yet the institution it defines persists. Equilibrium accounts are best understood as accounts of how institutions come into existence and are maintained in active use, and they need to be supplemented by an account of persistence through disuse.

\section{A model of institutional dynamics}

The persistence condition can be stated as a property of a simple dynamical model. Let the status of an institution at a time be one of four values: \emph{proposed}, \emph{active} (collectively accepted), \emph{dormant} (neither accepted nor rejected, but anchored), and \emph{ended} (rejected, or no longer anchored). Transitions between statuses are brought about by changes in collective attitudes.

\begin{definition}\label{d:transitions}
The \emph{admissible transitions} are: from proposed to active (by acceptance); from active to dormant (by lapse of acceptance); from dormant to active (by renewed acceptance); and from active or dormant to ended (by rejection or loss of anchoring). No transition leads out of ended, and no transition leads from proposed to dormant.
\end{definition}

\begin{proposition}\label{p:dynamics}
An institution exists at a time if and only if its status is active or dormant. Consequently, an institution that exists at one time continues to exist at every later time until a transition to ended occurs.
\end{proposition}

\begin{proof}
By Lemma~\ref{l:persist}, an institution that has come into existence persists until it is rejected or loses its anchoring, and coming into existence requires acceptance. The statuses active and dormant are exactly those that can be reached from proposed by acceptance without an intervening transition to ended. By Definition~\ref{d:transitions}, only rejection or loss of anchoring leads from these statuses to ended.
\end{proof}

The model makes explicit an asymmetry to which we return in replying to the objections below: entry into existence requires a positive attitude, exit requires a negative one, and the dormant status lies between them. The acceptance condition, by contrast, corresponds to a model with no dormant status, in which every lapse of acceptance is an exit. The cases of Section~3.2 are cases of dormancy, and it is because the acceptance model has no place for them that it misclassifies them.

\section{Recent work}

Three recent approaches bear on the argument. \citet{Epstein2015} distinguishes the grounding conditions of social facts, which say what makes a particular fact obtain, from their anchoring conditions, which say what puts those grounding conditions in place. Lemma~\ref{l:persist} is a claim about anchoring: it says that once the grounding conditions of an institution have been put in place, they remain in place until they are rejected by the anchoring community. Epstein's framework thus accommodates the lemma, and the cases of Section~3.2 provide evidence for a particular view of how anchoring works over time.

\citet{Asta2018} argues that many social properties are conferred: they are attributed to individuals by others who have standing to do so, on the basis of perceived base properties. Conferralism is an account of the instantiation of social properties by individuals rather than of the persistence of institutions, but it has a parallel consequence. A conferred property persists as long as the conferral is not withdrawn by those with standing, even if no one is currently attending to it. \citet{Haslanger2012} defends the reality of socially constructed kinds against the charge that their construction makes them less real, on grounds close to those of Corollary~\ref{c:real}. \citet{Mason2016} surveys the ways in which the metaphysics of social kinds has been integrated into general metaphysics, and argues that the same criteria of kindhood apply to social and to natural kinds, which is the assumption underlying our use of Principle~\ref{p:criterion}.

\section{Consequences}

The corrected persistence condition has consequences for several neighbouring questions. For political philosophy, it implies that states and their institutions do not lapse when they cease to command active allegiance, but only when they are rejected, which bears on the stability of political orders discussed in POL-07. For the philosophy of language, it suggests that linguistic conventions persist through periods of disuse, so that an archaic word retains its meaning until it is reinterpreted rather than merely forgotten, which is consistent with the treatment of meaning in LNG-01. For metaphysics, it locates social kinds alongside dispositions (MET-11): what persists through disuse is a disposition of the anchoring community not to reject, and this disposition is as real as the dispositions of physical objects.

\section{Objections}

\begin{objection}
Social kinds are mind-dependent, and whatever is mind-dependent is not part of objective reality.
\end{objection}

\begin{reply}
The objection equivocates on \enquote{mind-dependent}. Social kinds depend on minds for their existence, but facts about them are not up to any individual mind. Whether a coin is legal tender is not settled by what I believe about it, and I can be wrong about it. Objectivity in the relevant sense is independence from the beliefs of any particular judge, and social facts have it \citep{Thomasson2003}.
\end{reply}

\begin{objection}
The account is circular. Collective acceptance requires that members treat things as having a status, and treating something as money presupposes the concept of money.
\end{objection}

\begin{reply}
The concept of money is required, but not the existence of money. The members of a community can have the concept and act on it before any money exists, which is how money comes into existence. There is no circularity in the claim that a kind is brought into existence by attitudes whose content involves the concept of that kind.
\end{reply}

\begin{objection}
Lemma~\ref{l:persist} makes social kinds too persistent. On the account, an institution that everyone ignores persists indefinitely.
\end{objection}

\begin{reply}
An institution that everyone ignores persists as long as its rules remain anchored in the practices of the community. This is the correct result in the cases of Section~3.2. When the anchoring disappears, as when a society vanishes and its records are lost, the institution ceases to exist even without rejection. The account therefore does not imply indefinite persistence, only persistence through periods of disuse.
\end{reply}

\begin{objection}
The persistence condition overgenerates. An institution that was proposed but never accepted has never been rejected either, and so would persist.
\end{objection}

\begin{reply}
Lemma~\ref{l:persist} applies only to institutions that have come into existence, and coming into existence requires collective acceptance. A proposed institution that was never accepted never existed, and there is nothing for the absence of rejection to sustain. The asymmetry between creation and persistence is deliberate: it takes positive collective attitudes to bring an institution into being, but only their negative counterpart to end it.
\end{reply}

\begin{objection}
The account relies on the notion of the community whose practices contain an institution's rules, and this notion is indeterminate. Which community anchors a border: the citizens of the two states, their governments, or the international community?
\end{objection}

\begin{reply}
The notion is indeterminate in some cases, and in those cases the persistence of the institution is correspondingly indeterminate. This is a feature of the subject matter, not a defect of the account. Disputed borders are precisely the cases in which different communities reject and fail to reject the same status function, and in which there is no settled fact about whether the border persists. In the central cases, the anchoring community is fixed by the institution's own rules: a statute is anchored in the legal system that enacted it, a currency in the jurisdiction that issued it.
\end{reply}

\section{Conclusion}

Social kinds exist. They figure in explanations that cannot do without them, and their reality is not diminished by their dependence on collective attitudes. Their instances come into existence through collective acceptance but persist, once established, as long as they are not collectively rejected by a community whose practices still contain their rules. The persistence condition explains why dormant laws, unused currencies and vacant offices continue to exist, and why abolition, repeal and demonetization are the acts that end institutions.

\begin{thebibliography}{99}\small
\bibitem[{\'A}sta(2018)]{Asta2018} {\'A}sta (2018). \emph{Categories We Live By: The Construction of Sex, Gender, Race, and Other Social Categories}. Oxford University Press.
\bibitem[Durkheim(1895)]{Durkheim1895} Durkheim, {\'E}. (1895). \emph{Les R{\`e}gles de la m{\'e}thode sociologique}. Alcan.
\bibitem[Epstein(2015)]{Epstein2015} Epstein, B. (2015). \emph{The Ant Trap: Rebuilding the Foundations of the Social Sciences}. Oxford University Press.
\bibitem[Gilbert(1989)]{Gilbert1989} Gilbert, M. (1989). \emph{On Social Facts}. Routledge.
\bibitem[Guala(2016)]{Guala2016} Guala, F. (2016). \emph{Understanding Institutions: The Science and Philosophy of Living Together}. Princeton University Press.
\bibitem[Guala and Hindriks(2015)]{GualaHindriks2015} Guala, F. and Hindriks, F. (2015). A unified social ontology. \emph{The Philosophical Quarterly}, 65(259), 177--201.
\bibitem[Haslanger(2012)]{Haslanger2012} Haslanger, S. (2012). \emph{Resisting Reality: Social Construction and Social Critique}. Oxford University Press.
\bibitem[Lewis(1969)]{Lewis1969} Lewis, D. (1969). \emph{Convention: A Philosophical Study}. Harvard University Press.
\bibitem[Mason(2016)]{Mason2016} Mason, R. (2016). The metaphysics of social kinds. \emph{Philosophy Compass}, 11(12), 841--850.
\bibitem[Searle(1995)]{Searle1995} Searle, J.~R. (1995). \emph{The Construction of Social Reality}. Free Press.
\bibitem[Searle(2010)]{Searle2010} Searle, J.~R. (2010). \emph{Making the Social World: The Structure of Human Civilization}. Oxford University Press.
\bibitem[Thomasson(2003)]{Thomasson2003} Thomasson, A.~L. (2003). Foundations for a social ontology. \emph{ProtoSociology}, 18--19, 269--290.
\bibitem[Tuomela(2013)]{Tuomela2013} Tuomela, R. (2013). \emph{Social Ontology: Collective Intentionality and Group Agents}. Oxford University Press.
\end{thebibliography}

\end{document}
